% from ega2-1-fr.tex
\oldpage[II]{182}
$S=\bigoplus_{n\geqslant 0}\Gamma(X,\mathcal{L}^{\otimes n})$~;
l'hypothèse que $X$ est intègre entraîne que l'anneau $S$ est intègre
(\hyperref[I.7.4.4-fr]{I, 7.4.4}), donc il en est de même de $C$
(\hyperref[I.5.1.4-fr]{I, 5.1.4}). Pour démontrer que $C$ est normal, nous
utiliserons le lemme suivant~:

\begin{lemma}[8.8.6.1]
\label{II.8.8.6.1-fr}
Soit $Z$ un préschéma intègre normal. Alors l'anneau
$\Gamma(Z,\mathcal{O}_Z)$ est intègre et intégralement clos.
\end{lemma}

\begin{proof}
En effet, il résulte de
(\hyperref[I.8.2.1.1-fr]{I, 8.2.1.1}) que
$\Gamma(Z,\mathcal{O}_Z)$ est l'intersection, dans le corps des fonctions
rationnelles $R(Z)$, des anneaux intégralement clos $\mathcal{O}_z$ pour
$z\in Z$.
\end{proof}

Cela étant, montrons d'abord que $V$ est \emph{localement intègre et normal}~;
on peut pour cela se borner au cas où $X=\operatorname{Spec}(A)$ est affine,
d'anneau $A$ intègre et intégralement clos
(\hyperref[II.6.3.8-fr]{6.3.8}), et où
$\mathcal{L}=\mathcal{O}_X$. Comme alors
$V=\operatorname{Spec}(A[T])$ et que $A[T]$ est intègre et intégralement clos
([24], p.~99), cela établit notre assertion. D'autre part, pour tout ouvert
affine $U$ de $C$, $g^{-1}(U)$ est quasi-compact puisque le morphisme $g$ est
quasi-compact~; puisque $V$ est localement intègre, les composantes connexes
de $g^{-1}(U)$ sont des préschémas intègres ouverts dans $g^{-1}(U)$, donc en
nombre fini, et comme $V$ est normal, ces préschémas sont normaux
(\hyperref[II.6.3.8-fr]{6.3.8}). Alors $\Gamma(U,\mathcal{O}_C)$, qui est égal
à $\Gamma(g^{-1}(U),\mathcal{O}_V)$ en vertu de (i), est composé direct d'un
nombre fini d'anneaux intègres et intégralement clos
(\hyperref[II.8.8.6.1-fr]{8.8.6.1}), ce qui montre que $C$ est normal
(\hyperref[II.6.3.4-fr]{6.3.4}).
\end{proof}

\subsection{Le critère d'amplitude de Grauert : énoncé}
\label{subsection:II.8.9-fr}

Nous nous proposons de montrer que les propriétés démontrées dans
(\hyperref[II.8.8.2-fr]{8.8.2}) \emph{caractérisent} les
$\mathcal{O}_X$-Modules amples pour $f$, et de façon précise, de démontrer le
critère suivant~:

\begin{theorem}[8.9.1]
\label{II.8.9.1-fr}
\emph{(critère de Grauert).} Soient $Y$ un préschéma,
$p:X\to Y$ un morphisme séparé et quasi-compact, $\mathcal{L}$ un
$\mathcal{O}_X$-Module inversible. Pour que $\mathcal{L}$ soit ample
relativement à $p$, il faut et il suffit qu'il existe un $Y$-préschéma $C$,
une $Y$-section $\varepsilon:Y\to C$ de $C$, et un $Y$-morphisme
$q:\mathbf{V}(\mathcal{L})\to C$, ayant les propriétés suivantes~:
\begin{enumerate}
  \item[\rm{(i)}] Le diagramme
  \begin{equation}
  \label{II.8.9.1.1-fr}
    \xymatrix{
      X\ar[r]^{j}\ar[d]_{p} & \mathbf{V}(\mathcal{L})\ar[d]^{q}\\
      Y\ar[r]_{\varepsilon} & C
    }
  \tag{8.9.1.1}
  \end{equation}
  où $j$ est la section nulle du fibré vectoriel $\mathbf{V}(\mathcal{L})$,
  est commutatif.
  \item[\rm{(ii)}] La restriction de $q$ à
  $\mathbf{V}(\mathcal{L})-j(X)$ est une immersion ouverte quasi-compacte
  \[
    \mathbf{V}(\mathcal{L})-j(X)\longrightarrow C
  \]
  dont l'image ne rencontre pas $\varepsilon(Y)$.
\end{enumerate}
\end{theorem}

On notera que si $C$ est \emph{séparé} sur $Y$, on peut, dans la condition
(ii), supprimer l'hypothèse que l'immersion ouverte est quasi-compacte~; pour
voir en effet que cette dernière propriété est conséquence des autres
conditions, on peut se borner au cas où $Y$ est affine, et l'assertion découle
alors de (\hyperref[I.5.5.1-fr]{I, 5.5.1, (i)} et
\hyperref[I.5.5.10-fr]{5.5.10}). On peut aussi supprimer

